\documentclass[11pt]{article}

\usepackage[margin=1in]{geometry}
\usepackage{amsmath, amssymb, amsfonts}
\usepackage{booktabs}
\usepackage{longtable}
\usepackage{array}
\usepackage{enumitem}
\usepackage{hyperref}
\usepackage{xcolor}
\usepackage{listings}
\usepackage{caption}

\hypersetup{
    colorlinks=true,
    linkcolor=blue,
    citecolor=blue,
    urlcolor=blue
}

\lstset{
    basicstyle=\ttfamily\small,
    breaklines=true,
    frame=single,
    columns=fullflexible
}

\title{Wine Reviews Rerun Plan for FT+PPI Mean Estimation}
\author{Prepared for Internal Discussion}
\date{\today}

\begin{document}
\maketitle

\section{Objective of the Rerun}

This document provides a complete experimental redesign for rerunning the Wine Reviews mean-estimation experiment in the FT+PPI paper. The goal is not to make Wine Reviews the final main empirical application. Rather, the goal is to use Wine Reviews as a controlled benchmark to validate the full experimental pipeline before migrating the same workflow to more management-relevant mean-estimation tasks and to general M-estimation tasks.

The rerun should satisfy four requirements:

\begin{enumerate}[label=(\roman*)]
    \item All label accounting must be clean. If a label is used for ramp-up, validation, early stopping, hyperparameter selection, fine-tuning, or scaling-law fitting, this cost must be charged against the labeled budget.
    \item The ramp-up procedure must be part of the deployable estimator, not an oracle preprocessing step using extra labels.
    \item The experiment must include stronger inferential baselines, especially PPI++, Cross-PPI, calibrated PPI, and the old Qwen3-Embedding-0.6B + MLP replication baseline.
    \item The experiment must be implemented in a modular way so that the same codebase can later be reused for categorical response estimation and general M-estimation.
\end{enumerate}

The key point is that the Wine rerun should be treated as an implementation blueprint. The scientific standard should be high enough that, once the same logic is moved to a better dataset or M-estimation task, the empirical section can be defended without redesigning the methodology.

\section{Original Wine Experiment and Main Weaknesses}

The previous Wine Reviews experiment used the Wine Enthusiast dataset. After removing duplicated reviews, the empirical analysis contained roughly 119,955 unique reviews. The input was the review text only, and the outcome was the expert rating on an 80--100 scale. The old experiment treated 109,955 reviews as the target population and held out 10,000 additional labeled reviews to estimate the fine-tuning scaling law. The target population mean was approximately 88.441.

The old implementation used Qwen3-Embedding-0.6B as a text encoder. The final hidden-layer embedding was passed into a small multilayer perceptron regression head. Most encoder parameters were frozen; only the final transformer layer and the MLP head were updated. The ramp-up procedure was described in the online appendix: the labeled sample was split into a fine-tuning subset, a validation subset, and a PPI subset. The validation subset was used to estimate the scaling law and was not used for either fine-tuning or final PPI correction.

The old setup was useful as proof of concept, but it has several limitations that must be addressed in the rerun:

\begin{enumerate}[label=(\alph*)]
    \item The main scaling-law validation used an extra 10,000 labeled observations. This makes label accounting vulnerable under the limited-human-data framing.
    \item The ramp-up procedure was placed in the appendix rather than being implemented as part of the main estimator comparison.
    \item The main figures focused on estimated means and confidence intervals, which made the optimal allocation visually difficult to identify.
    \item The baselines were too limited. In particular, PPI++, Cross-PPI, calibrated PPI, and a clean replication of the old Qwen3-Embedding-0.6B + MLP pipeline were missing.
    \item The model architecture looked more like embedding regression than LLM fine-tuning. A rerun should use LoRA or QLoRA on an instruction-tuned causal language model.
\end{enumerate}

\section{Overall Structure of the New Wine Empirical Section}

The rerun should be organized as follows.

\begin{longtable}{p{0.22\textwidth}p{0.36\textwidth}p{0.34\textwidth}}
\toprule
Section & Purpose & Main Outputs \\
\midrule
Empirical setup and label accounting & Define data, estimand, labeled budget, unlabeled pool, and models & Data flow diagram and label accounting table \\
Ramp-up and scaling-law validation & Move ramp-up into the main experiment and test the residual-variance learning curve & Scaling-law curve, leave-one-size-out error, bootstrap intervals \\
Optimal allocation validation & Compare theoretical allocation, ramp-up allocation, and empirical oracle allocation & Allocation curve, normalized variance curve, regret table \\
Estimator comparison & Compare sample mean, PPI, PPI++, calibrated PPI, Qwen3-Embedding-0.6B + MLP, LoRA FT+PPI, and Cross-PPI & RMSE, bias, empirical SD, estimated SE, CI length, coverage, sample savings \\
Diagnostics and robustness & Address collapse, MSE versus variance loss, model scale, Cross-PPI, and external pilot labels & Residual variance, prediction variance, calibration, ablations \\
\bottomrule
\end{longtable}

\section{Data Setup}

\subsection{Cleaning Rules}

Use the same core setting as the original paper:
\[
X_i = \text{wine review description}, \qquad Y_i = \text{rating points} \in \{80,81,\ldots,100\}.
\]

The cleaning protocol should be fixed before running any experiment:

\begin{enumerate}
    \item Drop observations with missing \texttt{description} or missing \texttt{points}.
    \item Deduplicate by \texttt{description}. If the same description maps to multiple ratings, either retain the first occurrence or use a deterministic rule such as the modal rating. The chosen rule must be documented and fixed.
    \item Exclude all structured metadata from the main experiment, including variety, country, price, province, winery, and designation. This maintains comparability with the original experiment and prevents gains from being attributed to metadata rather than textual prediction.
    \item Record the finite population size $N$, finite population mean $\mu_N$, finite population variance $\sigma_Y^2$, the rating histogram, and the token-length distribution under each tokenizer.
\end{enumerate}

The rating distribution should be reported. Since the ratings are concentrated between 80 and 100 and often around the high 80s, this task may naturally make different methods look similar. This is not a reason to hide the limitation. It should be reported as part of the empirical context.

\subsection{Finite Population Protocol}

Treat the cleaned Wine Reviews dataset as a finite target population:
\[
\mathcal{P}=\{(X_i,Y_i)\}_{i=1}^N, \qquad
\mu_N = \frac{1}{N}\sum_{i=1}^N Y_i.
\]

For each replication $r$ and labeled budget $B$:

\begin{enumerate}
    \item Draw a labeled set $L_r \subset \mathcal{P}$ of size $B$ by simple random sampling without replacement.
    \item Define the unlabeled set as $U_r=\mathcal{P}\setminus L_r$. The experiment may use $X_i$ for $i\in U_r$ but may not use $Y_i$ for model training, model selection, scaling-law fitting, or estimator construction.
    \item All methods in the same replication must use the same $L_r$ and $U_r$. This enables paired method comparisons.
\end{enumerate}

Use the following labeled budgets:
\[
B\in\{1000,3000,5000,10000\}.
\]

A smaller setting such as $B=500$ can be used for debugging or appendix robustness, but it should not be part of the main table unless the ramp-up cost is handled carefully.

Recommended replication counts:

\begin{itemize}
    \item Debug stage: $R=3$.
    \item Pilot stage: $R=10$.
    \item Full Wine rerun: target $R=30$; if GPU resources are constrained, use at least $R=20$ for the main tables.
\end{itemize}

\section{Model Architecture}

\subsection{Main LoRA / QLoRA Model}

The rerun should replace the old embedding-plus-MLP main model with LoRA or QLoRA fine-tuning of an instruction-tuned causal language model.

Recommended model tiers:

\begin{longtable}{p{0.22\textwidth}p{0.32\textwidth}p{0.36\textwidth}}
\toprule
Tier & Model & Purpose \\
\midrule
Debug & Qwen2.5-1.5B-Instruct or another small Qwen model & Fast pipeline check for losses, PPI estimators, and Hyak jobs \\
Main & Qwen2.5-7B-Instruct & Main Wine rerun model \\
Robustness & Llama-3.1-8B-Instruct or Qwen3-4B/8B & Architecture or model-family robustness \\
\bottomrule
\end{longtable}

The main model should treat the wine rating as a continuous scalar outcome, not as a discrete classification problem. Use a regression head on top of the instruction-tuned model representation, following the original continuous-regression setup but replacing the old embedding encoder with a LoRA/QLoRA-adapted causal language model.

For each input review, use a fixed text template such as:

\begin{lstlisting}
System: You are an expert wine critic. Given a wine review, predict the expert rating on an 80-100 scale.

User:
Review:
{description}

Target:
Continuous rating score on the 80-100 scale.
\end{lstlisting}

Let $h_\theta(X_i)$ denote the pooled final hidden representation of this prompted input, for example the final-token representation or another fixed pooling rule chosen before the experiment. The surrogate prediction is:
\[
f_\theta(X_i)=w^\top h_\theta(X_i)+b,
\qquad f_\theta(X_i)\in\mathbb{R}.
\]

The model is trained against the observed numerical rating $Y_i$ as a continuous response. Do not normalize over the discrete labels $\{80,\ldots,100\}$, and do not train the main Wine rerun as a classification task.

This design has three advantages. First, it preserves the mean-estimation setting. Second, it stays close to the original Wine experiment's continuous regression target, making the rerun easier to compare with the current paper. Third, it keeps the prediction object as a scalar $f_\theta(X)$, which is exactly the quantity needed by PPI, PPI++, and Cross-PPI.

\subsection{Embedding Baseline}

The rerun should retain only the old embedding-regression pipeline as the embedding / ML baseline:

\begin{longtable}{p{0.28\textwidth}p{0.28\textwidth}p{0.32\textwidth}}
\toprule
Baseline & Feature Representation & Learner \\
\midrule
Qwen3-Embedding-0.6B + MLP & Frozen or partially tuned embedding & Old-paper replication baseline \\
\bottomrule
\end{longtable}

Do not add other classical ML or embedding baselines to the Wine rerun. Keeping only Qwen3-Embedding-0.6B + MLP makes the comparison focused: the rerun asks whether the new label accounting, ramp-up, and PPI variants improve the old Wine pipeline.

\section{LoRA / QLoRA Training Configuration}

\subsection{Default Configuration}

Recommended default settings:

\begin{longtable}{p{0.32\textwidth}p{0.52\textwidth}}
\toprule
Parameter & Recommended Value \\
\midrule
Quantization & 4-bit NF4, bfloat16 compute \\
LoRA rank $r$ & 16 \\
LoRA alpha & 32 \\
LoRA dropout & 0.05 \\
Target modules & \texttt{q\_proj}, \texttt{k\_proj}, \texttt{v\_proj}, \texttt{o\_proj}, \texttt{gate\_proj}, \texttt{up\_proj}, \texttt{down\_proj} \\
Max sequence length & 256 first; if more than 5\% of examples are truncated, use 384 \\
Optimizer & paged\_adamw\_8bit or AdamW \\
Learning rate & $2\times 10^{-4}$ \\
Learning-rate robustness & $1\times 10^{-4}$ and $5\times 10^{-5}$ \\
Weight decay & 0 \\
Warmup ratio & 0.03 \\
Scheduler & cosine \\
Epochs & 3 for $s\geq 1000$; 5 for $s<1000$ \\
Early stopping & Patience of 2 evaluations \\
Evaluation frequency & Every 0.25 epoch or every 50 steps \\
Effective batch size & At least 32 if possible; 64 preferred \\
\bottomrule
\end{longtable}

\subsection{Important Implementation Issue: Variance Loss and Micro-Batches}

The residual-variance loss is a batch-level objective:
\[
L_{\mathrm{var}}
=
\frac{1}{q}\sum_{i=1}^{q}(r_i-\bar r)^2,
\qquad
r_i=Y_i-f_\theta(X_i),
\qquad
\bar r=\frac{1}{q}\sum_{i=1}^{q} r_i.
\]

If the micro-batch size is too small, the batch variance estimate will be noisy and unstable. Standard gradient accumulation is not automatically equivalent to computing the variance over the full effective batch, because each micro-batch may already backpropagate a different batch-centered loss.

There are two acceptable implementations:

\begin{enumerate}
    \item Use a true forward batch size of at least 32 whenever GPU memory allows.
    \item Implement custom accumulated loss: forward multiple micro-batches, collect residuals, compute the variance loss over the concatenated effective batch, and only then backpropagate.
\end{enumerate}

A rating-stratified batch sampler is recommended so that each batch contains observations from multiple rating levels. This reduces instability when ratings are concentrated near the population mean.

\subsection{Training Losses}

The LoRA pipeline should implement two primary continuous-regression losses.

\paragraph{MSE Loss.}
\[
L_{\mathrm{MSE}}=\frac{1}{q}\sum_{i=1}^{q}\left(Y_i-f_\theta(X_i)\right)^2.
\]
This is included to align with the original MSE baseline.

\paragraph{Residual-Variance Loss.}
\[
L_{\mathrm{Var}}
=
\frac{1}{q}\sum_{i=1}^{q}
\left[
\left(Y_i-f_\theta(X_i)\right)
-
\frac{1}{q}\sum_{j=1}^{q}\left(Y_j-f_\theta(X_j)\right)
\right]^2.
\]
This is the main proposed fine-tuning objective.

As a robustness check, one may also evaluate an anchored version:
\[
L_{\mathrm{AnchoredVar}}=L_{\mathrm{Var}}+\lambda \bar r^2,
\qquad \lambda\in\{0.01,0.05\}.
\]
This should not be the main method because it changes the theoretical objective. It is only a diagnostic for identifiability and collapse concerns.

\section{Ramp-Up Procedure with Clean Label Accounting}

The ramp-up procedure should be moved into the main empirical design. The main result should use target-only label accounting. The old external 10,000-label procedure can be reported only as a separate robustness regime.

\subsection{Main Regime: Target-Only Conservative Ramp-Up}

Given a labeled budget $B$, draw a validation set:
\[
V_r\subset L_r, \qquad |V_r|=v_B.
\]

Recommended validation sizes:

\begin{longtable}{ccc}
\toprule
Budget $B$ & Validation size $v_B$ & Effective final budget $n_{\mathrm{eff}}=B-v_B$ \\
\midrule
1000 & 200 & 800 \\
3000 & 300 & 2700 \\
5000 & 500 & 4500 \\
10000 & 1000 & 9000 \\
\bottomrule
\end{longtable}

The remaining labeled set is:
\[
L_r'=L_r\setminus V_r, \qquad n_{\mathrm{eff}}=B-v_B.
\]

Construct a nested sequence of fine-tuning sets:
\[
T_{s_1}\subset T_{s_2}\subset \cdots \subset T_{s_K}\subset L_r',
\qquad |T_{s_\ell}|=s_\ell.
\]

Recommended ramp-up grids:

\begin{longtable}{ccc}
\toprule
Budget $B$ & Validation size $v_B$ & Fine-tuning size grid $s_\ell$ \\
\midrule
1000 & 200 & 50, 100, 150, 250, 400 \\
3000 & 300 & 100, 200, 400, 700, 1000, 1500 \\
5000 & 500 & 100, 200, 400, 700, 1000, 1500, 2500 \\
10000 & 1000 & 100, 200, 400, 700, 1000, 1500, 2500, 4000 \\
\bottomrule
\end{longtable}

For each $s_\ell$:

\begin{enumerate}
    \item Fine-tune a LoRA model on $T_{s_\ell}$.
    \item Evaluate validation residual variance on $V_r$:
    \[
    \widehat{\nu}(s_\ell)=\widehat{\operatorname{Var}}_{V_r}\{Y-f^{(s_\ell)}(X)\}.
    \]
    \item Fit the residual-variance scaling law:
    \[
    \widehat{\nu}(s)=\hat a s^{-\hat\alpha}+\hat b.
    \]
    \item Choose the allocation by solving:
    \[
    \hat s^*
    =
    \arg\min_{0<s<n_{\mathrm{eff}}}
    \frac{\hat a s^{-\hat\alpha}+\hat b}{n_{\mathrm{eff}}-s}.
    \]
    \item Choose the nearest feasible grid point or continue ramp-up if the predicted optimum exceeds the current largest training size.
    \item Train the final model on $T_{\hat s}$ and use the correction set
    \[
    C_r=L_r'\setminus T_{\hat s}
    \]
    for PPI correction.
\end{enumerate}

The validation set $V_r$ must not be used in final PPI correction, because it has been used for scaling-law fitting, checkpoint selection, and model selection. This conservative protocol charges the validation labels to the method. It is stricter than the original paper and therefore more defensible.

\subsection{Robustness Regime: External Pilot Labels Available}

As a separate practical regime, replicate the spirit of the original experiment:

\begin{enumerate}
    \item Set aside 10,000 labeled reviews as an external pilot set $H$.
    \item Use $H$ to estimate scaling-law parameters, tune the prompt format, and choose hyperparameters.
    \item Use the remaining reviews as the target population.
    \item For each target labeled budget $B$, do not charge $H$ against the current-task budget, but clearly report this as an external-pilot setting.
\end{enumerate}

This should not be the main result. It is a practical scenario corresponding to firms that have historical expert ratings or survey archives.

\subsection{Optional Robustness: No-Waste Cross-Validated Ramp-Up}

If compute resources allow, add an appendix version that does not permanently set aside a validation set. For each candidate fine-tuning size $s$, estimate residual variance by cross-validation inside $T_s$, then train the final model on all of $T_s$ and reserve $L_r\setminus T_s$ for PPI correction.

This version saves labels but is much more expensive because it requires multiple fine-tuning runs per candidate size. It should be optional.

\section{Scaling-Law Validation}

The scaling-law section should no longer rely only on in-sample $R^2$. It should validate both the curve and the allocation decision.

\subsection{Scaling Curve}

For $B=10000$, plot:
\[
s\mapsto \widehat{\nu}(s)=\widehat{\operatorname{Var}}_{V_r}\{Y-f^{(s)}(X)\}.
\]

The figure should include:

\begin{enumerate}
    \item The fitted power-law curve in linear space.
    \item A log-log plot of $\widehat\nu(s)-\hat b$ against $s$.
    \item Bootstrap bands across replications.
    \item A constant-predictor baseline $\widehat{\operatorname{Var}}(Y)$.
\end{enumerate}

The main panel can focus on LoRA-Var, while appendix panels show LoRA-MSE.

\subsection{Leave-One-Size-Out Prediction Error}

For each fine-tuning size $s_\ell$ in the grid:

\begin{enumerate}
    \item Remove the point $(s_\ell,\widehat\nu(s_\ell))$.
    \item Fit $\hat a$, $\hat\alpha$, and $\hat b$ using the remaining sizes.
    \item Predict the held-out residual variance.
    \item Report the normalized prediction error:
    \[
    \mathrm{LOSO\ Error}
    =
    \frac{|\widehat\nu(s_\ell)-\widehat\nu_{-\ell}(s_\ell)|}
    {\widehat\nu(s_\ell)}.
    \]
\end{enumerate}

This directly addresses the concern that a smooth parametric curve may overfit noisy small-sample observations.

\subsection{Allocation Sensitivity}

Bootstrap the fitted scaling-law parameters and report:

\begin{itemize}
    \item confidence intervals for $\hat a$, $\hat\alpha$, and $\hat b$;
    \item confidence intervals for $\hat s^*/B$;
    \item allocation regret:
    \[
    \mathrm{Regret}(\hat s)
    =
    \frac{V(\hat s)-V(s^{\mathrm{oracle}})}{V(s^{\mathrm{oracle}})}.
    \]
\end{itemize}

The allocation decision matters more than the exact parameter estimates. If $\hat s^*$ is not exact but regret is low, the method is still practically useful.

\section{Optimal Allocation Validation}

The original allocation figure should be replaced. Instead of plotting estimated means and confidence intervals, plot normalized estimator variance or normalized CI length directly.

\subsection{Oracle Allocation Grid}

For $B=10000$, evaluate the following allocation ratios:
\[
\frac{s}{B}\in
\{0,0.025,0.05,0.075,0.10,0.125,0.15,0.20,0.30,0.40,0.50,0.70,0.90\}.
\]

If $s<50$, treat the method as no-FT or PPI-only rather than training an unstable model.

For each allocation $s$:

\begin{enumerate}
    \item Train $f^{(s)}$ using the $s$ fine-tuning labels.
    \item Use $B-s$ labels for PPI correction.
    \item Use all unlabeled observations in $U_r$ for the prediction term.
    \item Repeat over replications and compute estimator distribution, empirical variance, confidence interval length, and coverage.
\end{enumerate}

This oracle grid is not a deployable method. It is used to validate the theory and to show the shape of the allocation tradeoff.

\subsection{Recommended Figure}

The new allocation figure should have three panels.

\paragraph{Panel A: Residual-Variance Learning Curve.}
\[
\widehat\nu(s)=\widehat{\operatorname{Var}}\{Y-f^{(s)}(X)\}.
\]

\paragraph{Panel B: Normalized Estimator Variance or CI Width.}
Estimate:
\[
\widehat V_{\mathrm{PPI}}(s)
=
\frac{\widehat{\nu}_C(s)}{B-s}
+
\frac{\widehat\tau_U(s)}{|U_r|},
\qquad
\widehat\tau_U(s)=\widehat{\operatorname{Var}}_{U_r}\{f^{(s)}(X)\}.
\]

Plot:
\[
\frac{\widehat V_{\mathrm{PPI}}(s)}{\widehat V_{\mathrm{SampleMean}}}.
\]

This is much clearer than plotting estimated means.

\paragraph{Panel C: Efficiency Gain.}
\[
R^2(s)-\frac{s}{B},
\qquad
R^2(s)=1-\frac{\widehat\nu(s)}{\widehat\sigma_Y^2}.
\]

This panel directly corresponds to the theoretical criterion for outperforming the sample mean.

\subsection{Allocation Validation Table}

Report:

\begin{longtable}{p{0.22\textwidth}ccccc}
\toprule
Method & $\hat s^*/B$ & $s^{\mathrm{oracle}}/B$ & Absolute gap & Regret & Normalized CI length \\
\midrule
LoRA-Var &  &  &  &  &  \\
LoRA-MSE &  &  &  &  &  \\
\bottomrule
\end{longtable}

MSE does not have a theoretically privileged downstream loss, but its residual-variance curve can still be fitted and used to compute an allocation benchmark.

\section{Estimator Comparison}

The main comparison should be organized by estimator family. All comparisons should use paired replications.

\subsection{Label-Only Baseline}

\paragraph{Sample Mean.}
\[
\hat\mu_{\mathrm{mean}}=\frac{1}{B}\sum_{i\in L_r}Y_i.
\]
This uses all $B$ labels and is the most important baseline.

\subsection{Fixed-Predictor Rectification Baselines}

\paragraph{Zero-Shot LLM + PPI.}
Use the base LLM without any label-based fine-tuning. All $B$ labels are used for correction.

\paragraph{Zero-Shot LLM + PPI++ / Linear Calibrated PPI.}
For mean estimation, use the control-variate form:
\[
\hat\mu_{\lambda}
=
\bar Y_C+\hat\lambda(\bar f_U-\bar f_C).
\]
A finite-sample estimate is:
\[
\hat\lambda
=
\frac{\widehat{\operatorname{Cov}}_C(Y,f)}{\widehat{\operatorname{Var}}_C(f)}
\cdot
\frac{|U|}{|U|+|C|}.
\]
To avoid overfitting $\lambda$, use a 5-fold cross-fitted $\lambda$: for each correction fold, estimate $\lambda$ using the other folds.

\subsection{Qwen3 Embedding Replication Baseline}

Include:

\begin{enumerate}
    \item Qwen3-Embedding-0.6B + MLP + PPI.
    \item Qwen3-Embedding-0.6B + MLP + PPI++.
\end{enumerate}

For fair comparison, report two allocation versions when feasible:

\begin{itemize}
    \item same allocation as the LoRA-Var ramp-up rule;
    \item own oracle allocation over the grid, reported as an upper-bound baseline.
\end{itemize}

\subsection{Fine-Tuning Only}

Report LoRA-FT-only methods using MSE and Var losses. These use all $B$ labels for fine-tuning and estimate the mean by averaging $f(X)$ over the population. These estimators are generally biased. Therefore, do not report them as valid inferential methods with PPI confidence intervals. Report RMSE and bias only.

\subsection{Split FT+PPI}

Main methods:

\begin{enumerate}
    \item LoRA-MSE + PPI.
    \item LoRA-Var + PPI.
\end{enumerate}

Each should be reported in two versions:

\begin{enumerate}
    \item Oracle allocation: grid-search best allocation, reported as an upper bound.
    \item Ramp-up allocation: deployable main result with all ramp-up labels counted.
\end{enumerate}

\subsection{Split FT+PPI++}

Add:

\begin{enumerate}
    \item LoRA-MSE + PPI++.
    \item LoRA-Var + PPI++.
\end{enumerate}

This is important. If LoRA-Var still improves after PPI++ calibration, then the contribution is not merely due to using a weak rectification baseline.

\subsection{Cross-PPI}

Use $K=5$ folds. Split $L_r$ into $I_1,\ldots,I_K$. For fold $k$, train $f_{-k}$ without using the labels in $I_k$. Estimate:
\[
\hat\mu_{\mathrm{CrossPPI}}
=
\frac{1}{K}\sum_{k=1}^{K}\frac{1}{|U_r|}\sum_{u\in U_r} f_{-k}(X_u)
+
\frac{1}{B}\sum_{k=1}^{K}\sum_{i\in I_k}\left(Y_i-f_{-k}(X_i)\right).
\]

Run at least:

\begin{enumerate}
    \item Cross-PPI-MSE.
    \item Cross-PPI-Var.
\end{enumerate}

Cross-PPI is not just another baseline. It directly addresses the question of whether sample splitting is necessary. The strongest result would be that Cross-PPI-Var improves over Cross-PPI-MSE, showing that the variance-aligned objective complements cross-fitting.

\subsection{Main Comparison Table}

For each budget $B$, report:

{\small
\setlength{\tabcolsep}{3pt}
\begin{longtable}{@{}p{0.25\textwidth}ccccccc@{}}
\toprule
Method & RMSE & Bias & Emp. SD & Est. SE & CI Length & Coverage & Sample Savings \\
\midrule
Sample Mean &  &  &  &  &  &  &  \\
PPI &  &  &  &  &  &  &  \\
PPI++ &  &  &  &  &  &  &  \\
Qwen3-Embedding-0.6B + MLP + PPI++ &  &  &  &  &  &  &  \\
LoRA-MSE + PPI++ &  &  &  &  &  &  &  \\
LoRA-Var + PPI++ &  &  &  &  &  &  &  \\
Cross-PPI-MSE &  &  &  &  &  &  &  \\
Cross-PPI-Var &  &  &  &  &  &  &  \\
\bottomrule
\end{longtable}
}

Definitions:
\[
\mathrm{RMSE}=\sqrt{\frac{1}{R}\sum_{r=1}^{R}(\hat\mu_r-\mu_N)^2},
\]
\[
\mathrm{Bias}=\frac{1}{R}\sum_{r=1}^{R}(\hat\mu_r-\mu_N),
\]
\[
\mathrm{SampleSavings}
=
1-
\frac{\widehat{\operatorname{Var}}(\hat\mu_{\mathrm{method}})}
{\widehat{\operatorname{Var}}(\hat\mu_{\mathrm{SampleMean}})}.
\]

All method comparisons should report paired standard errors. For a metric $M$, define:
\[
d_r=M_{r,A}-M_{r,B}.
\]
Report $\bar d$ and its paired standard error rather than only reporting unpaired standard errors for each estimator.

\section{Confidence Intervals and Variance Estimation}

\subsection{Vanilla PPI}

For a fixed predictor $f$:
\[
\hat\mu_{\mathrm{PPI}}=\bar f_U+\overline{Y-f(X)}_C.
\]

Estimated variance:
\[
\widehat{\operatorname{Var}}(\hat\mu)
=
\frac{\widehat{\operatorname{Var}}_C(Y-f(X))}{|C|}
+
\frac{\widehat{\operatorname{Var}}_U(f(X))}{|U|}.
\]

The 95\% confidence interval is:
\[
\hat\mu \pm 1.96\sqrt{\widehat{\operatorname{Var}}(\hat\mu)}.
\]

\subsection{PPI++ / Calibrated PPI}

For:
\[
\hat\mu_\lambda=\bar Y_C+\hat\lambda(\bar f_U-\bar f_C),
\]
use:
\[
\widehat{\operatorname{Var}}(\hat\mu_\lambda)
=
\frac{\widehat{\operatorname{Var}}_C(Y-\hat\lambda f(X))}{|C|}
+
\frac{\hat\lambda^2\widehat{\operatorname{Var}}_U(f(X))}{|U|}.
\]

If $\hat\lambda$ is cross-fitted, compute residuals fold by fold using the corresponding out-of-fold $\hat\lambda$ and then aggregate the variance.

\subsection{Coverage}

Coverage is evaluated against the finite population mean $\mu_N$:
\[
\mathbf{1}\{\mu_N\in CI_r\}.
\]

The main text can use the standard PPI variance estimator. A finite-population correction can be included as an appendix robustness check.

\section{Diagnostics and Robustness}

Every trained model should save the following diagnostics on validation, correction, and unlabeled sets whenever labels are available.

\subsection{Prediction Diagnostics}

Report:

\begin{itemize}
    \item residual mean $\overline{Y-f(X)}$;
    \item residual variance $\widehat{\operatorname{Var}}(Y-f(X))$;
    \item MSE;
    \item MAE;
    \item prediction variance $\widehat{\operatorname{Var}}(f(X))$;
    \item Pearson and Spearman correlations between $Y$ and $f(X)$;
    \item calibration intercept and slope from $Y=\alpha+\beta f(X)+\varepsilon$.
\end{itemize}

\subsection{Collapse Diagnostics}

To address representation-collapse concerns, report:

\begin{enumerate}
    \item whether $\widehat{\operatorname{Var}}(f(X))$ is close to zero;
    \item whether residual variance is lower than the constant-predictor baseline;
    \item the histogram and quantiles of continuous predictions $f(X)$;
    \item the share of predictions outside the observed rating range $[80,100]$ before any optional clipping.
\end{enumerate}

If Var loss outputs nearly constant predictions without reducing residual variance, that is a failure mode. If Var loss preserves nontrivial prediction variance and lowers residual variance, it directly addresses the collapse criticism.

\subsection{MSE versus Var Mechanism Figure}

Add an appendix figure with:

\begin{itemize}
    \item $x$-axis: fine-tuning size $s$;
    \item $y$-axis: validation residual variance;
    \item curves: MSE, Var, and constant predictor.
\end{itemize}

A second panel can show validation MSE. The goal is to show that MSE may perform well on prediction error but does not necessarily minimize the residual variance that controls PPI efficiency.

\section{Hyak Execution Plan}

UW Hyak uses Slurm, so the code should be designed around Slurm job arrays and checkpointed training.

\subsection{Recommended Code Structure}

\begin{lstlisting}
src/
  data/
    build_wine.py
    split_protocol.py
  models/
    regression_model.py
    lora_trainer.py
    embedding_baselines.py
  losses/
    mse_regression_loss.py
    residual_variance_loss.py
  estimators/
    mean.py
    ppi.py
    ppi_plus.py
    cross_ppi.py
  allocation/
    scaling_fit.py
    ramp_up.py
    oracle_grid.py
  experiments/
    exp00_debug.py
    exp01_scaling_law.py
    exp02_allocation.py
    exp03_estimator_comparison.py
    exp04_cross_ppi.py
  analysis/
    make_tables.py
    make_figures.py
configs/
  wine_debug.yaml
  wine_full_qwen25_7b.yaml
  wine_crossppi.yaml
slurm/
  run_lora_array.sbatch
  run_eval_array.sbatch
\end{lstlisting}

Training and estimation should be separated. Each trained model should produce predictions. PPI, PPI++, and Cross-PPI should be recomputed offline from saved prediction files.

\subsection{Prediction Cache Schema}

Save predictions in a long-format Parquet file with columns such as:

\begin{lstlisting}
run_id
replication_id
budget_B
method
loss
model_name
allocation_s
fold_id
sample_id
split_role       # train / validation / correction / unlabeled
y_true           # missing or masked for unlabeled in estimator code
pred_mean
pred_raw
pred_var_optional
pred_clipped_optional
\end{lstlisting}

This schema is designed for the Wine continuous-regression rerun. It can later be extended with probability columns if a separate categorical response task is added.

\subsection{Slurm Job Array Design}

Each LoRA training run should correspond to a Slurm array task:

\begin{lstlisting}
TASK = (model, loss, budget, replication, allocation_s, fold_id)
\end{lstlisting}

For split FT+PPI, use \texttt{fold\_id = 0}. For Cross-PPI, use \texttt{fold\_id = 1, ..., K}.

\section{Execution Phases}

\subsection{Phase 0: Non-LLM Unit Tests}

Use synthetic data to test:

\begin{itemize}
    \item PPI unbiasedness;
    \item PPI++ variance behavior;
    \item Cross-PPI fold logic;
    \item confidence interval coverage under simple data-generating processes.
\end{itemize}

No GPU is needed.

\subsection{Phase 1: Wine Embedding Baseline}

Run:

\begin{itemize}
    \item sample mean;
    \item Qwen3-Embedding-0.6B + MLP old-paper replication.
\end{itemize}

The purpose is to validate the split protocol, label accounting, estimator code, and table generation.

\subsection{Phase 2: LoRA Debug}

Use Qwen2.5-1.5B-Instruct.

\begin{itemize}
    \item Budgets: $B=1000$ and $B=5000$.
    \item Replications: $R=3$.
    \item Losses: MSE, Var.
    \item Allocations: $s/B\in\{0.1,0.2\}$.
\end{itemize}

Check training stability, residual variance, prediction variance, PPI coverage, GPU memory, and runtime.

\subsection{Phase 3: Scaling Law and Allocation}

Use Qwen2.5-7B-Instruct.

\begin{itemize}
    \item Budget: $B=10000$.
    \item Replications: start with $R=10$, then increase to $R=20$ or $R=30$ if stable.
    \item Allocation grid:
    \[
    s/B\in\{0.025,0.05,0.075,0.10,0.125,0.15,0.20,0.30,0.40,0.50\}.
    \]
\end{itemize}

Run larger allocation ratios such as 0.70 and 0.90 only if needed to display the full U-shape.

\subsection{Phase 4: Estimator Comparison}

Use Qwen2.5-7B-Instruct.

\begin{itemize}
    \item Budgets: $B\in\{1000,3000,5000,10000\}$.
    \item Replications: at least $R=20$, target $R=30$.
    \item Methods: sample mean, PPI, PPI++, Qwen3-Embedding-0.6B + MLP + PPI/PPI++, LoRA-MSE + PPI/PPI++, LoRA-Var + PPI/PPI++, ramp-up allocation, oracle allocation upper bounds.
\end{itemize}

\subsection{Phase 5: Cross-PPI}

Run Cross-PPI after the split-pipeline is stable.

\begin{itemize}
    \item Folds: $K=5$.
    \item Budgets: $B=1000$ and $B=5000$ first; add $B=10000$ if resources allow.
    \item Replications: start with $R=10$.
    \item Losses: MSE, Var.
\end{itemize}

Cross-PPI is about five times more expensive than a single split run, so it should not be run at full scale before the rest of the pipeline is stable.

\section{Draft Structure for the New Empirical Section}

\subsection{Empirical Setup}

Emphasize that Wine Reviews is a mean-estimation benchmark. Do not oversell it as a full human-behavior simulation task. The task is estimating average expert ratings from textual expert evaluations.

\subsection{Data-Driven Ramp-Up and Scaling-Law Diagnostics}

State that the ramp-up procedure is implemented as part of the estimator, not as an oracle preprocessing step. Validation labels are counted. The section should report in-sample fit, leave-one-size-out prediction error, and allocation regret.

\subsection{Optimal Allocation Validation}

Show the U-shaped allocation curve using normalized estimator variance or normalized CI width. Report the gap between the scaling-law allocation and empirical oracle allocation.

\subsection{Comparison with Baselines}

Compare against PPI++, Cross-PPI, the Qwen3-Embedding-0.6B + MLP replication baseline, and standard continuous fine-tuning losses. Report coverage and CI length in addition to RMSE.

\subsection{Robustness}

Include collapse diagnostics, external-pilot labels, model-size robustness, anchored variance loss, and Cross-PPI with variance loss.

\section{Design Errors to Avoid}

\begin{enumerate}
    \item Do not use an extra 10,000 labels in the main method while sample mean does not receive them.
    \item Do not use validation labels for scaling-law fitting and then reuse them for PPI correction.
    \item Do not only plot estimated means against allocation ratios. Plot estimator variance, CI length, and normalized variance ratio.
    \item Do not compare only against vanilla PPI. PPI++ and Cross-PPI must appear.
    \item Do not report only unpaired standard errors. Use paired comparisons across replications.
    \item Do not compute variance loss on tiny micro-batches unless a correct accumulated-loss implementation is used.
    \item Do not present FT-only as a valid inferential estimator. It may be reported for RMSE and bias, but not for valid PPI confidence intervals.
\end{enumerate}

\section{Minimum Viable Rerun and Full Rerun}

\subsection{Minimum Viable Rerun}

Used to confirm the Hyak pipeline and method implementation:

\begin{itemize}
    \item Model: Qwen2.5-1.5B-Instruct.
    \item Budgets: $B=1000$ and $B=5000$.
    \item Replications: $R=3$.
    \item Losses: MSE, Var.
    \item Methods: sample mean, PPI, PPI++, split FT+PPI, Cross-PPI.
    \item Allocations: $s/B\in\{0.1,0.2\}$.
    \item Diagnostics: residual variance, prediction variance, RMSE, coverage.
\end{itemize}

\subsection{Full Wine Rerun}

Used for a serious appendix or internal validation:

\begin{itemize}
    \item Model: Qwen2.5-7B-Instruct.
    \item Budgets: $B\in\{1000,3000,5000,10000\}$.
    \item Replications: at least $R=20$, target $R=30$.
    \item Losses: MSE, Var.
    \item Allocation validation: $B=10000$ grid over $s/B$.
    \item Main estimator comparison: ramp-up selected allocation, oracle allocation upper bounds, PPI and PPI++.
    \item Cross-PPI: $B=1000$ and $B=5000$, $K=5$, $R=10$; add $B=10000$ if resources allow.
    \item Baselines: sample mean, PPI-only, PPI++, old Qwen3-Embedding-0.6B + MLP replication.
\end{itemize}

\section{Codex Implementation Prompt}

The following prompt can be given directly to Codex.

\begin{lstlisting}
Implement the Wine Reviews rerun pipeline for FT+PPI mean estimation.

Requirements:
1. Clean Wine Reviews data and construct a finite population protocol.
2. Implement reproducible labeled/unlabeled splits for budgets B in {1000,3000,5000,10000}.
3. Implement sample mean, vanilla PPI, PPI++ with cross-fitted lambda, and Cross-PPI estimators.
4. Implement continuous regression prediction for ratings 80-100 using a scalar regression head on the LoRA/QLoRA model representation.
5. Implement LoRA/QLoRA fine-tuning with MSE loss and residual-variance loss.
6. Implement conservative ramp-up with validation labels counted and excluded from PPI correction.
7. Implement scaling-law fitting and allocation selection.
8. Implement oracle allocation grid for B=10000.
9. Implement the frozen embedding baseline: Qwen3-Embedding-0.6B+MLP.
10. Save all predictions and estimator outputs in parquet/JSON so estimators can be recomputed without retraining.
11. Add leakage checks: no correction label may be used in training, validation, scaling-law fitting, or checkpoint selection for the same split model.
12. Produce tables for RMSE, bias, empirical SD, estimated SE, CI length, coverage, sample savings, and paired differences.
13. Produce figures for scaling law, allocation variance curve, efficiency gain R^2(s)-s/B, and collapse diagnostics.
\end{lstlisting}

Minimum leakage checks:

\begin{lstlisting}
assert train_ids.isdisjoint(correction_ids)
assert validation_ids.isdisjoint(correction_ids)
assert prompt_demo_ids.isdisjoint(correction_ids)  # if prompt baselines are used
# For Cross-PPI fold k:
assert fold_k_ids.isdisjoint(train_ids_for_f_minus_k)
\end{lstlisting}

\section{Success Criteria}

The Wine rerun should be considered successful if it establishes the following points:

\begin{enumerate}
    \item The residual-variance scaling law is empirically stable, not only in terms of $R^2$ but also in leave-one-size-out prediction and allocation regret.
    \item The ramp-up selected allocation is close to the empirical oracle allocation, or at least has low regret.
    \item The variance loss reduces residual variance relative to MSE.
    \item The advantage of variance loss remains after PPI++ or calibrated PPI is applied.
    \item Cross-PPI improves label reuse, but Cross-PPI-Var further improves over Cross-PPI-MSE, showing that variance-aligned fine-tuning and cross-fitting are complementary.
    \item Confidence interval coverage remains valid.
    \item Collapse diagnostics show nontrivial prediction variance and lower residual variance than the constant-predictor baseline.
\end{enumerate}

If these conditions are met on Wine Reviews, the experimental pipeline is ready to be transferred to survey mean estimation, categorical response estimation, and M-estimation tasks.

\end{document}
